%
% Documento: Introdução
%

\chapter{Introdução}\label{chap:introducao}
Este template em \LaTeX \ foi criado para o programa de Pós-Graduação em Modelagem Computacional em Ciência e Tecnologia da Universidade Estadual de Santa Cruz (UESC), mas também pode ser utilizado por outros programas de graduação e pós-graduação da UESC e de outras instituições, desde que sejam seguidas as respectivas regras de formatação.

A classe {\ttfamily ppgmc-uesc.cls} foi elaborada com base nas normas da ABNT e da UESC para trabalhos de conclusão de curso, e maiores detalhes relacionados aos comandos existentes no estilo podem ser adquiridos através da documentação disponível no site \href{http://www.abntex.net.br/}{http://www.abntex.net.br/}. Encorajamos a consultar as regras do pacote abntex  \href{http://linorg.usp.br/CTAN/macros/latex/contrib/abntex2/doc/abntex2.pdf}{neste link}.

Para melhor entendimento do uso do estilo de formatação, recomendamos a análise dos comandos existentes no arquivo {\ttfamily PPGMC.tex} e os resultados obtidos após a compilação. Os arquivos pré-textuais, textuais e pós-textuais podem ser modificados conforme a necessidade, e as referências citadas são apenas exemplos.

A utilização deste template requer conhecimento prévio de \LaTeX, e sugestões de melhorias e correções são sempre bem-vindas. O template está disponível para uso online na plataforma \href{https://pt.overleaf.com}{\textit{Overleaf}}, e na seção \ref{sec:contatos} estão os contatos do autor para dúvidas e contribuições.

\section{Motivação}
\label{sec:motivacao}

Este template foi desenvolvido para padronizar os trabalhos de conclusão de curso do programa de Pós-Graduação em Modelagem Computacional em Ciência e Tecnologia da UESC e demais programas de mestrado e doutorado da universidade. O estilo de documento utilizado é o {\ttfamily abntex2}, que possui comandos especiais para auxiliar a distribuição/definição das diversas partes constituintes do projeto.

Uma das principais vantagens do uso deste template é a formatação \textit{automática} dos elementos que compõem um documento acadêmico, tais como capa, folha de rosto, dedicatória, agradecimentos, resumo, abstract, listas de figuras, tabelas, siglas e símbolos, sumário, capítulos e referências. Isso torna a elaboração do trabalho mais ágil e eficiente, permitindo que o autor possa focar na elaboração do conteúdo.
